\section{The formal language for reading statements: EasyCrypt}
\label{sec:formal-language}

\subsection{Core predicates and relations}

\begin{definition}[Packed-key prefix]
\label{def:vk-prefix}
\[
  \operatorname{VKPrefix}_{n}(sk,vk)
  \;\Longleftrightarrow\;
  \forall i\;(0\le i<n\Rightarrow sk[i]=vk[i]).
\]
This is \identifier{vk_prefix_eq} in the code; mode~2 uses \(n=992\).
\end{definition}

\begin{definition}[Prefix agreement of external key regions]
\label{def:external-prefix}
\[
  \operatorname{ExternalKeyPrefix}(mem,sku,vku)
  \Longleftrightarrow
  \forall i\in[0,992),\;
  \operatorname{load8}(mem,sku+i)=\operatorname{load8}(mem,vku+i).
\]
This is \identifier{external_key_prefix_match}. It compares the first
992 bytes exported by actual raw KeyGen to distinct external SK/VK addresses.
\end{definition}

\begin{definition}[Prefix in memory]
\label{def:mem-prefix}
\[
  \operatorname{SKMemPrefix}(sk,mem,b)
  \;\Longleftrightarrow\;
  \forall i\in[0,992),\;
  sk[i]=\operatorname{load8}(mem,\operatorname{uint}(b)+i).
\]
This predicate compares Sign's local \identifier{BArray2752} with Verify's
raw-memory reads in one relation.
\end{definition}

\begin{definition}[Canonical address range]
\label{def:canonical-region}
For an integer address \(a\) and length \(n\),
\[
\begin{aligned}
  \operatorname{CanonicalAddress}(a)
    &\Longleftrightarrow 0\le a<2^{64},\\
  \operatorname{CanonicalRegion}(a,n)
    &\Longleftrightarrow
      0\le a\;\land\;0\le n\;\land\;a+n\le 2^{64}.
\end{aligned}
\]
The first premise permits a raw application binary interface (raw ABI)
integer address to round-trip through \identifier{W64.of_int} and
\identifier{W64.to_uint}; the second excludes modular wraparound of the address
range. For length 0, the second formula alone also permits \(a=2^{64}\), so the
actual round-trip theorems separately require \identifier{canonical_ui64_address}
for the necessary addresses and lengths. This is a call contract, not a proof
of physical memory safety.
\end{definition}

\begin{definition}[Memory-region equality]
\label{def:region-eq}
\[
  \operatorname{RegionEq}(m_1,m_2,b,n)
  \Longleftrightarrow
  \forall i\in[0,\operatorname{uint}(n)),\;
  \operatorname{load8}(m_1,\operatorname{uint}(b)+i)
  =\operatorname{load8}(m_2,\operatorname{uint}(b)+i).
\]
This is \identifier{region_eq} in the code. Week~6 hash theorems require only
this relation on the regions actually read in VK, pre, and message, rather than
whole-memory equality.
\end{definition}

\begin{definition}[Canonical signature-prefix inputs]
\label{def:signature-prefix-canonical}
\identifier{canonical_challenge} means that each of the 256 challenge words
is exactly 0 or 1. \identifier{canonical_signed_low} means that each of the
1024 low words equals the sign extension of its low byte, with signed integer
value in \([-128,128)\). The actual prefix round trip holds under these two premises.
\end{definition}

\begin{definition}[Canonical mode-2 HBZ and leaf frames]
\label{def:canonical-hbz}
For a 32-bit coefficient array \(h\) of capacity \(2048\),
\[
  \operatorname{CanonicalHBZ}_2(h)
  \Longleftrightarrow
  \forall i\in[0,1024),\;-6\le\operatorname{sint}(h[i])<7.
\]
This is \identifier{canonical_hbz_mode2} in the code.
\identifier{prepared_hbz_prefix} says that the first \(n\) symbol bytes equal
\(\operatorname{sint}(h[i])+6\); \identifier{decoded_hbz_prefix} says that the
first \(n\) recovered coefficients equal the originals.
\identifier{byte_tail_frame} and \identifier{coeff_tail_frame} record that
\([1024,2048)\) is unchanged before and after leaf execution in their respective arrays.
\end{definition}

\begin{definition}[Mode-2 HBZ table certificate]
\label{def:hbz-table-certificate}
\identifier{mode2_hbz_table_certificate} is the conjunction of three facts.
\begin{enumerate}
  \item The four fields of the actual encoder table for each of the 13 symbols equal their specifications.
  \item Actual decoder words contain each interval's start/frequency exactly.
  \item The actual 512 packed symbol words map each of 1024 slots to its unique interval symbol.
\end{enumerate}
\identifier{actual_mode2_hbz_tables_certified} applies this predicate to the
literal arrays in the pack/unpack extractions. It does not assume correctness
of the external certificate-generation script: the kernel rechecks every
generated EasyCrypt equality.
\end{definition}

\begin{definition}[Canonical rANS byte trace]
\label{def:rans-byte-trace}
For a canonical symbol list \(s=(s_0,\ldots,s_{1023})\),
\identifier{encode_trace} processes symbols in reverse from \(x_{1024}=2^{23}\),
determining boundary states \((x_j)_{j=0}^{1024}\) and decoder-order byte
segments \((\beta_j)_{j=0}^{1023}\). The cursor cuts and complete byte sequence are
\[
  c_0=4,\qquad c_{j+1}=c_j+|\beta_j|,
  \qquad
  B=\operatorname{serialize32\_le}(x_0)
      \Vert\beta_0\Vert\cdots\Vert\beta_{1023}
\]
The code's \identifier{valid_rans_trace}(s,xs,cuts,bytes) is simply the
conjunction equating its arguments to these three objects computed by
\identifier{trace_states}, \identifier{trace_cuts}, and
\identifier{trace_bytes}. It assumes neither existence of an external trace
nor decoder success; \identifier{valid_rans_trace_canonical} proves the
canonical construction.
\end{definition}

\begin{definition}[Array/list bridge and segment frame]
\label{def:rans-array-list-bridge}
For a byte array \(a\) of capacity 2048, \identifier{symbol_list_of_array}(a)
is the integer list of its first 1024 elements, and
\identifier{symbol_suffix}(a,i) drops the first \(i\) elements from that list.
\identifier{segment_matches}(a,p,bs) means
\[
  0\le p\;\land\;p+|bs|\le2048\;\land\;
  \forall k\in[0,|bs|),\quad
  \operatorname{uint}(a[p+k])=bs[k]
\]
The predicate \identifier{prefix_frame}(before,after,p) says that the arrays'
\(k\)-th bytes agree for every \(k<p\). Week~10's compiled theorems provide
\[
\begin{aligned}
  \mathsf{symbol\_suffix}(a,1024)&=[],\\
  \mathsf{symbol\_suffix}(a,i)
    &=\operatorname{uint}(a[i])::\mathsf{symbol\_suffix}(a,i+1)
\end{aligned}
\]
together with rules for segment concatenation, prepending a one-byte write,
and shrinking frames.
\end{definition}

\begin{definition}[Actual encoder inner-step invariant]
\label{def:rans-actual-inner}
\identifier{encoder_inner_phase_inv} and
\identifier{encoder_inner_segment_inv} directly open the generated
\identifier{_rans_encode} inner normalization loop. They assert existence of
entry state \((x_0,off_0)\) and step \(k\) such that
\[
\begin{gathered}
  4\le off_0\le1024,\qquad 0\le k\le4,
  \qquad 1\le\operatorname{uint}(x_{\max}),\\
  \operatorname{uint}(off)=off_0-k,\qquad
  \operatorname{uint}(x)=\mathsf{encoder\_phase\_state}(x_0,k),\\
  \mathsf{segment\_matches}
    (encp,off_0-k,\mathsf{encoder\_phase\_written}(x_0,k))
\end{gathered}
\]
and preservation of the lower prefix relative to the array at entry to that
inner call. These are Week~10's conservative local predicates, verifying the
actual one-byte store and W64 cursor safety. Week~11's stronger tail-aware
predicate adds outer accumulation and the exact 0/1/2-byte length.
\end{definition}

\begin{definition}[Tail-aware inner invariant]
\label{def:rans-tail-aware-inner}
Week~11's \identifier{encoder_inner_tail_inv} combines the exact pure
normalization length \(\ell=\mathsf{mode2\_normalization\_len}(x_0,s)\) with the
encoding trace of the existing pure symbol tail \(tail\). Its core segment condition is
\[
  \mathsf{segment\_matches}\bigl(
    encp,off_0-k,
    \mathsf{inner\_written\_suffix}(x_0,s,k)\mathbin{++}
      \mathsf{encode\_trace}(tail).\mathsf{bytes}
  \bigr)
\]
along with \(0\le k\le\ell\), cursor and state equalities, and the prefix
frame. The predicate is preserved through actual-guard progress and exit,
yielding exactly \(k=\ell\le2\). It is then used in the outer one-symbol
recurrence and final-store composition to close the actual outer procedure's
complete successful postcondition.
\end{definition}

\begin{verified}[Success-conditioned complete suffix of the actual encoder]
\label{thm:rans-encoder-suffix}
\identifier{actual_rans_encode_trace_closure} directly opens the actual
generated \identifier{_rans_encode}, proving either returned \(bad\ne0\) or,
when \(bad=0\),
\[
  \mathsf{segment\_matches}
  (encp',\operatorname{uint}(off'),
   \mathsf{trace\_bytes}(\mathsf{symbol\_list\_of\_array}(symbols_0)))
\]
together with \(off'+|\mathsf{trace\_bytes}|=1024\), the successful size range,
and \identifier{prefix_frame}. Success is not a precondition. Because this is
Hoare partial correctness, termination and success reachability remain separate obligations.
\end{verified}

\begin{definition}[Actual decoder exact-trace input]
\label{def:rans-decoder-exact-input}
\identifier{actual_mode2_decoder_trace_input}(a,b,st,n) is a functional
representation contract between expected symbol array \(a\), encoded buffer
\(b\), initial state \(st\), and logical size \(n\). Its core conditions are
\[
\begin{gathered}
  \mathsf{mode2\_hbz\_symbol\_stream}(a),\qquad
  n=|\mathsf{trace\_bytes}(\mathsf{symbol\_list\_of\_array}(a))|,\\
  4\le n\le1024,\qquad
  \mathsf{segment\_matches}(b,0,\mathsf{trace\_bytes}(\cdots)),\\
  (st[0],st[1],st[2])=(1024,n,13)
\end{gathered}
\]
and pointwise read equalities for every normalization segment. The actual
mode-2 \identifier{symbol_words} and \identifier{dsyms_words} arrays are fixed
by separate bindings in the top-level Hoare theorem. This predicate does not
assume \(bad=0\), an output-array equality, or a post-state.
\end{definition}

\begin{definition}[Decoder outer/inner trace invariants]
\label{def:rans-decoder-trace-invariants}
At iteration \(i\), \identifier{decoder_outer_trace_live} binds the actual
state and cursor to
\[
  x=\mathsf{decoder\_state\_at}(s,i),\qquad
  \operatorname{uint}(off)=\mathsf{decoder\_cursor}(s,i)
\]
while maintaining the expected symbols in \(decoded[0..i)\) and the input-array
frame in \(decoded[i..2048)\). After a one-symbol word update,
\identifier{decoder_inner_trace_live} uses
\[
\begin{aligned}
  k&=\operatorname{uint}(off)-\mathsf{decoder\_cursor}(s,i),\\
  x&=\mathsf{decoder\_replay\_prefix}
       (\mathsf{decoder\_state\_at}(s,i+1),s_i,k)
\end{aligned}
\]
and maintains \(0\le k\le|\mathsf{decoder\_segment\_at}(s,i)|\). The generated
guard agrees exactly with \(k<|\mathsf{segment}|\), and the segment length is
0, 1, or 2 bytes. Thus each actual read is either a one-byte consumption
transition or a no-consumption exit to the next outer boundary.
\end{definition}

\begin{verified}[Actual decoder exact-trace recovery]
\label{thm:rans-decoder-trace-refinement}
\identifier{actual_rans_decode_trace_refinement} directly opens actual
generated \identifier{RansDecodeTarget.M._rans_decode} under
\cref{def:rans-decoder-exact-input} and bindings for the two actual tables.
Every terminating execution satisfies
\[
  bad'=0,\qquad off'=n,\qquad
  decoded'[0..1024)=a[0..1024)
\]
and the output tail \([1024,2048)\) equals the initial array. The internal
\(x=2^{23}\) at outer-loop exit excludes the actual final rejection check,
but is absent from the ABI postcondition. This is exact-trace Hoare partial
correctness, not a proof of termination or actual encoder success.
\end{verified}

\begin{definition}[Actual rANS harness control]
\label{def:rans-actual-harness}
\identifier{Mode2RansActualHarness.run} calls actual \identifier{_rans_encode},
then calls actual \identifier{__copy_encoded_suffix} and
\identifier{_rans_decode} only when returned \(bad\) is 0. In that success
branch, decoder-state input fields are exactly
\((\mathsf{count},\mathsf{size},m)=(1024,\mathsf{encoded\_size},13)\).
The current harness-control theorem concludes only
\[
  bad_{\mathrm{enc}}\in\{0,1\},\qquad
  \mathsf{decoder\_ran}\Longleftrightarrow bad_{\mathrm{enc}}=0
\]
and this field initialization. That control theorem is not itself an inverse.
Week~13's separate semantic theorem sequentially invokes
\cref{thm:rans-encoder-suffix}, \cref{thm:rans-decoder-trace-refinement}, and the
actual-copy theorem inside this module to close the following conclusion.
\end{definition}

\begin{verified}[Actual encoder--copy--decoder core inverse]
\label{thm:rans-core-actual-inverse}
\identifier{actual_rans_encode_copy_decode_inverse} assumes only initial
argument bindings and \identifier{mode2_hbz_symbol_stream}. In a terminating
execution, actual encoder failure leaves the decoder unexecuted and preserves
its initial array/state. On actual encoder success, the copied buffer becomes
an exact \identifier{trace_bytes} input and the decoder satisfies
\[
  bad_{\mathrm{dec}}=0,\qquad
  off_{\mathrm{dec}}=\mathsf{encoded\_size},\qquad
  decoded[0..1024)=symbols[0..1024)
\]
while preserving the \([1024,2048)\) tail. Encoder success, copied-region
equality, decoder success, and output equality are not preconditions. This is
Hoare partial correctness; success reachability and termination of both loops
remain separate obligations.
\end{verified}

\begin{definition}[Production full-HBZ harness]
\label{def:production-full-hbz-harness}
\identifier{SignaturePackUnpackHbzFullHarness.run} calls actual production
\identifier{SignaturePackMode2Target.M._encode_hb_z1_full}, then calls actual
\identifier{SignatureUnpackMode2Target.M._decode_hb_z1_full} only when the
returned size is not 0. \identifier{HbzFullActualHarness.run} executes the
same control pattern with focused full-HBZ extractions. Week~14's relational
theorem exactly aligns the result tuples and \identifier{decoder_ran} flags
of the two harnesses.
\end{definition}

\begin{verified}[Production full-HBZ success-conditioned inverse]
\label{thm:production-full-hbz-inverse}
\identifier{signature_pack_unpack_hbz_full_actual_exact} proves that the
two harnesses return equal values and agree on decoder execution from the
same initial arguments. The subsequent
\identifier{signature_pack_unpack_hbz_full_inverse_mode2} assumes only initial
argument bindings and \identifier{canonical_hbz_mode2}. For a terminating
execution, \(size=0\) leaves the decoder unexecuted and preserves the initial
encoded/decoded/bad arrays. When \(size\ne0\),
\[
  4\le\operatorname{uint}(size)\le1024,
  \qquad bad=0,
  \qquad decoded[0..1024)=hbz_0[0..1024)
\]
and the coefficient tail, encoded-suffix frame, and actual
\identifier{prepared_symbols} trace witness are preserved. Encoder success,
decoder success, output equality, and the trace are not preconditions. This is
Hoare partial correctness; reaching nonzero size and termination are separate obligations.
\end{verified}

\begin{boundary}[Success-conditioned HBZ]
\label{bdry:hbz-success}
\(\operatorname{CanonicalHBZ}_2(h)\) guarantees success of the prepare leaf,
but not that the actual full encoder with its fixed 1024-byte rANS workspace
returns nonzero size. Week~14's production theorem does not assume encoder or
decoder success and retains the actual return's \(\mathsf{size}=0\) branch in
its conclusion. Week~15 excludes that failure branch only for terminating
fixed all-zero-input executions, not for general canonical inputs or termination itself.
\end{boundary}

\begin{verified}[Production success for fixed all-zero HBZ]
\label{thm:fixed-zero-hbz-success}
The only preconditions of
\identifier{signature_pack_unpack_hbz_zero_success_mode2} bind the initial
encoded/decoded/bad arrays and the exact \identifier{zero_hbz} input. A
terminating execution satisfies
\[
  size\ne0,\qquad 4\le\operatorname{uint}(size)\le1024,\qquad
  \mathsf{decoder\_ran},\qquad bad=0
\]
with all first 1024 recovered coefficients equal to 0, preserving the
coefficient/encoded tail frames and exact
\identifier{mode2_trace_bytes(all_six_symbols)}. The subtheorem
\identifier{actual_rans_encode_all_six_success} shows that nonzero \(bad\) from
the actual encoder requires a normalization trace exceeding 1020 bytes,
contradicting the all-6 trace's 1020-byte bound. This is fixed-input Hoare
partial correctness, not termination/losslessness or a \(\operatorname{phoare}=1\) theorem.
\end{verified}

\begin{definition}[Actual KeyGen matrix/finalizer snapshot]
\label{def:keygen-matrix-finalizer-snapshot}
\leavevmode\par
\identifier{ActualM23MatrixFinalizeSnapshot.run} stores the returned \(bp\)
from actual \identifier{_kp_m23_matrix} as \(\mathsf{pre\_bp}\), then calls
actual \identifier{_keypair_finalize_m23}. At every active coefficient,
\[
\begin{aligned}
  r_i&=(\mathsf{pre\_bp}_i+s_{2,i}+a_i)\bmod q,\\
  r_i&=2b_{1,i}+b_{0,i},\qquad -1\le b_{0,i}\le1,\\
  s'_{2,i}&=s_{2,i}-b_{0,i}
\end{aligned}
\]
defines \identifier{actual_snapshot_low_high_decomposition}.
Over the same active range, \identifier{actual_snapshot_mod2q_zero} requires
\[
  2(a_i-2b_{1,i})+
  2\mathsf{pre\_bp}_i+2s'_{2,i}\equiv0\pmod{2q}
\]
the displayed equality.
\end{definition}

\begin{verified}[Semantic composition of actual KeyGen snapshots]
\label{thm:keygen-semantic-snapshot}
\identifier{actual_m23_matrix_finalize_semantic_snapshot} assumes only exact
bindings of six initial arrays, matrix/input representation ranges, active
\(s_2\in[-1,1]\), and active \(a\in[0,q)\). Its conclusion combines the matrix
output/NTT representation, exact finalizer semantics, the two predicates of
\cref{def:keygen-matrix-finalizer-snapshot}, and initial--snapshot--final tail
frames. Neither finalizer output nor the KG equation is a precondition.
\end{verified}

\begin{boundary}[\identifier{STOP-KG-NTT} at paper freeze]
\label{bdry:keygen-ntt-mul}
The \(\mathsf{pre\_bp}\) in \cref{thm:keygen-semantic-snapshot} is actual
\identifier{output_row}. Target 77's
\identifier{output_row_from_mode2_ntt_words} expands it into the actual
full-invNTT/pointwise-word representation, and
\identifier{actual_m23_matrix_snapshot_rows_explicit} transports it to the two
actual harness rows. Both results compile, but do not identify that expression
with the security model's polynomial product
\(A_{\mathrm{gen}}s_{\mathrm{gen}}\). Odd-root orthogonality/full-NTT convolution
and the \identifier{Rq.poly}-to-security-list multiplication adapter are absent
from the frozen targets, so the snapshot congruence cannot be promoted to the
paper's \(As=qj\pmod{2q}\). This audit stopped the frozen paper's KeyGen frontier
at KG-2/finalization.
\end{boundary}

\begin{verified}[Post-freeze Rq--HAETAE representation preservation]
\label{thm:rq-haetae-post-freeze}
\identifier{RqHAETAEBridge.rq_poly_repr} requires only HAETAE polynomial
well-formedness and \identifier{incoeff} equality on the active 256
coefficients. For this relation,
\identifier{rq_poly_add_repr}, \identifier{rq_poly_mul_repr},
\identifier{rq_poly_dot3_repr} preserve addition, checked negacyclic
multiplication, and 3-term row products, respectively.
\identifier{Mode2KeygenNttMulBridge}'s \identifier{keygen_mode2_row_product} and
\identifier{mode2_row_product_haetae_dot} connect actual mode-2 NTT output to
this relation. The desired \(As=qj\) equation appears in none of the premises.
\end{verified}

\begin{boundary}[Latest \identifier{STOP-KG-AVEC-QJ}]
\label{bdry:keygen-avec-qj}
The post-freeze actual-sampler theorem interprets terminating
\identifier{avec} through SHAKE128 bytes with nonces 515/516, little-endian
16-bit candidates, and \(<64513\) rejection, connecting it to the paper's
\(a\). The paper's \(qj\) is a separate basis correction with \(q=64513\) only
at row 0, coefficient 0. In contrast, the current security model's named
\(qj\) is a synthetic seed generator, always giving 401 at that position.
\identifier{security_named_qj_is_not_paper_qj} compiles this inequality, so the
two objects cannot be identified. The first repair obligation is
\identifier{OBL-KG-SECURITY-EXPANDVECA-SEMANTICS}.
\end{boundary}

\begin{definition}[\(\mu\) prefix]
\label{def:mu-prefix}
Let \(\mu_S^{64}\) be Sign's 64-byte array and \(\mu_V^{32}\) Verify's
32-byte array. Then
\[
  \operatorname{MuPrefix}_{32}(\mu_S^{64},\mu_V^{32})
  \;\Longleftrightarrow\;
  \forall i\in[0,32),\;\mu_S^{64}[i]=\mu_V^{32}[i].
\]
that is, \(\take_{32}(\mu_S^{64})=\mu_V^{32}\).
\end{definition}

\begin{definition}[Raw branch and valid range]
\label{def:raw-region}
For 64-bit words \(p,\ell\), the two code predicates are
\[
  \operatorname{RawPreLen}(\ell)
    \Longleftrightarrow 0\le\operatorname{uint}(\ell)<2^{63},
  \qquad
  \operatorname{ValidRegion}(p,\ell)
    \Longleftrightarrow
    \operatorname{uint}(p)+\operatorname{uint}(\ell)\le 2^{64}
\]
The first selects the raw branch in Verify's actual branch on
\identifier{prelen >> 63}; the second excludes wraparound in address addition.
\end{definition}

\subsection{Hoare statements and relational equivalence}

The EasyCrypt statement
\[
  \operatorname{hoare}[P:\Phi\Longrightarrow\Psi]
\]
does not assert termination of \(P\). It states partial correctness: every
execution starting in \(\Phi\) that \textbf{terminates} satisfies \(\Psi\).
Losslessness or almost-sure termination of KeyGen's retry loop, and equality
with the paper's distribution, therefore require separate theorems.

The relational statement
\[
  \operatorname{equiv}[P\sim Q:\Phi\Longrightarrow\Psi]
\]
means that results of \(P,Q\), started in states satisfying the pre-relation
\(\Phi\), satisfy post-relation \(\Psi\). The hash theorems here do not replace
SHAKE with an abstract function; they directly compare the two sides' generated
absorb/finalize/Keccak/squeeze procedures in this relational logic.

\subsection{Exact traces and accepted tails}

The trace modules of the new raw-API theorems copy actual procedure control
flow while adding only ghost observations such as \identifier{observed_mu},
pointers, lengths, and \identifier{tail_reached}. Exact equivalence means
that actual procedure and trace yield the same final \identifier{Glob.mem}
and return value, not that every execution reaches the hash tail.

Verify can reject before passing signature-length, unpacking, norm, and
other checks. The current theorem exposes this branch by proving
\[
  \operatorname{VerifyAccept}\Longrightarrow\operatorname{TailReached}
\]
and only then, for successful executions, connects the VK/pre/message
addresses and lengths passed to actual \identifier{__verify_hash_mu} to the
observations. This is distinct from claiming that every signature reaches the tail.

Preservation of the actual procedure's result by an exact trace does not
automatically place the internal hash-call return value and post-execution
ghost fields in a relational theorem. The remaining
\claim{OBL-MU-PRODUCT-REPLAY} is the rule connecting
\identifier{observed_mu} in two completed sequential traces to the existing
pRHL hash-return theorem. Simple \identifier{sim}, \identifier{seq}, and
\identifier{byequiv} cannot eliminate the different residual continuations,
so this edge is explicitly \statusdeferred.

\subsection{Array capacity and logical length}

The extracted theories use fixed arrays capable of supporting several parameter sets:
\identifier{BArray2080}, \identifier{BArray2752},
\identifier{BArray2948}. This does not mean that mode-2 logical lengths
are 2080, 2752, and 2948 respectively. Actual mode-2 API lengths are 992 bytes
for the verification key, 1408 for the secret key, and 1474 for the signature;
the current theorems track the relation on the first 992 bytes. Separately,
signature-codec results track the layout and round trip of 32 challenge bytes
and 1024 low bytes, the first 1056 signature bytes. In HBZ leaves,
\identifier{BArray8192} holds 2048 32-bit coefficients and symbol
\identifier{BArray2048} holds 2048 bytes, but the mode-2 logical count is 1024.
Week~8 frame theorems therefore explicitly preserve the \([1024,2048)\) tails
of both arrays.
